\documentclass[aspectratio=169]{beamer}

% ─── Tema y colores ──────────────────────────────────────────────────────────
\usetheme{Madrid}
\usecolortheme{default}
\setbeamertemplate{navigation symbols}{}
\setbeamertemplate{itemize items}[circle]

% ─── Encoding y lenguaje ─────────────────────────────────────────────────────
\usepackage[utf8]{inputenc}
\usepackage[T1]{fontenc}
\usepackage[spanish,es-noquoting]{babel}

% ─── Gráficos y diagramas ────────────────────────────────────────────────────
\usepackage{tikz}
\usetikzlibrary{shapes.geometric,arrows.meta,positioning,fit,
                backgrounds,calc,shadows}

% ─── Código fuente ───────────────────────────────────────────────────────────
\usepackage{listings}
\usepackage{xcolor}

\definecolor{codegreen}{rgb}{0,0.52,0}
\definecolor{codegray}{rgb}{0.5,0.5,0.5}
\definecolor{codepurple}{rgb}{0.5,0,0.62}
\definecolor{codeblue}{rgb}{0.0,0.18,0.72}
\definecolor{backcolour}{rgb}{0.96,0.96,0.94}
\definecolor{codered}{rgb}{0.68,0.1,0.1}
\definecolor{keyorange}{rgb}{0.75,0.35,0.0}

\lstdefinestyle{pythonstyle}{
    backgroundcolor=\color{backcolour},
    commentstyle=\color{codegreen}\itshape,
    keywordstyle=\color{codeblue}\bfseries,
    numberstyle=\tiny\color{codegray},
    stringstyle=\color{codepurple},
    basicstyle=\ttfamily\fontsize{6.8}{8.2}\selectfont,
    breakatwhitespace=false,
    breaklines=true,
    keepspaces=true,
    numbers=left,
    numbersep=4pt,
    showstringspaces=false,
    showtabs=false,
    tabsize=2,
    language=Python,
    frame=single,
    rulecolor=\color{black!30},
    xleftmargin=10pt,
    morekeywords={from,import,as,def,return,True,False,None,
                  class,self,with,yield,lambda,async,await},
}

\lstdefinestyle{jsonstyle}{
    backgroundcolor=\color{backcolour},
    basicstyle=\ttfamily\fontsize{6.8}{8.2}\selectfont,
    breaklines=true,
    frame=single,
    rulecolor=\color{black!30},
    showstringspaces=false,
    xleftmargin=10pt,
    morestring=[b]",
    stringstyle=\color{codepurple},
}

\lstdefinestyle{shellstyle}{
    backgroundcolor=\color{gray!10},
    basicstyle=\ttfamily\fontsize{7}{8.5}\selectfont,
    breaklines=true,
    frame=single,
    rulecolor=\color{black!30},
    showstringspaces=false,
}

% ─── Otros paquetes ──────────────────────────────────────────────────────────
\usepackage{booktabs}
\usepackage{amsmath}
\usepackage{array}

% ─── Información del documento ───────────────────────────────────────────────
\title[IA Embebida en Sistemas]{\textbf{IA Embebida en Sistemas Computacionales}}
\subtitle{LLM \textbullet{} RAG \textbullet{} Redis \textbullet{} LangChain \textbullet{} Ollama}
\author[Ing. Informática]{\textbf{Curso: Ingeniería de Software con IA}}
\institute[Universidad]{Departamento de Ciencias de la Computación}
\date{Semestre I -- 2026}

% ─────────────────────────────────────────────────────────────────────────────
\begin{document}

% ═══════════════════════════════════════════════════════════════════
% DIAPOSITIVA 1: PORTADA
% ═══════════════════════════════════════════════════════════════════
\begin{frame}
  \titlepage
\end{frame}

% ═══════════════════════════════════════════════════════════════════
% DIAPOSITIVA 2: AGENDA
% ═══════════════════════════════════════════════════════════════════
\begin{frame}{Agenda de la Clase}
  \begin{columns}
    \column{0.5\textwidth}
      \begin{enumerate}
        \item Motivación y objetivos
        \item Problema de negocio
        \item Arquitectura tradicional \textit{vs.} moderna
        \item Componente RAG: fundamentos y pipeline
        \item Componente Redis: memoria contextual
        \item Componente LangChain: orquestación LCEL
      \end{enumerate}
    \column{0.5\textwidth}
      \begin{enumerate}
        \setcounter{enumi}{6}
        \item Pre-Prompt Engineering
        \item Salidas estructuradas en JSON
        \item Implementación paso a paso
        \item Actividades guiadas (3)
        \item Preguntas de reflexión
        \item Arquitectura final integrada
        \item Conclusiones y próxima clase
      \end{enumerate}
  \end{columns}
  \vspace{0.8em}
  \begin{exampleblock}{Prerrequisitos asumidos}
    Ollama local, embeddings, chunking básico, RAG introductorio, recuperación documental.
  \end{exampleblock}
\end{frame}

% ═══════════════════════════════════════════════════════════════════
% DIAPOSITIVA 3: OBJETIVOS
% ═══════════════════════════════════════════════════════════════════
\begin{frame}{Objetivos de Aprendizaje}
  \begin{columns}
    \column{0.5\textwidth}
    \begin{block}{Comprensión arquitectural}
      \begin{enumerate}
        \item Comprender un LLM como \textbf{componente} dentro de un sistema mayor, no como solución en sí mismo.
        \item Diseñar una arquitectura que integre Ollama, LangChain, Redis y RAG con responsabilidades bien definidas.
        \item Aplicar Prompt Engineering para restringir el comportamiento del LLM en un dominio específico.
      \end{enumerate}
    \end{block}
    \column{0.5\textwidth}
    \begin{exampleblock}{Habilidades prácticas}
      \begin{enumerate}
        \setcounter{enumi}{3}
        \item Implementar un pipeline RAG funcional con \texttt{ChromaDB} y \texttt{OllamaEmbeddings}.
        \item Utilizar Redis como memoria de sesión persistente mediante \texttt{RedisChatMessageHistory}.
        \item Construir una cadena LCEL que produzca objetos JSON validados con Pydantic, consumibles por otras aplicaciones.
      \end{enumerate}
    \end{exampleblock}
  \end{columns}
\end{frame}

% ═══════════════════════════════════════════════════════════════════
% DIAPOSITIVA 4: MOTIVACIÓN
% ═══════════════════════════════════════════════════════════════════
\begin{frame}{Motivación: Del LLM Aislado al Sistema Inteligente}
  \begin{columns}
    \column{0.6\textwidth}
    \begin{alertblock}{El error conceptual más común en proyectos de IA}
      Tratar el LLM como si fuera \textbf{la aplicación} en lugar de tratarlo
      como uno de sus componentes. Un LLM sin contexto, sin memoria y sin
      restricciones estructurales produce resultados impredecibles e
      inintegrables.
    \end{alertblock}
    \vspace{0.5em}
    \begin{block}{La narrativa de esta clase}
      \centering
      \large
      $\text{LLM solo}$
      $\;\longrightarrow\;$
      $\text{LLM + RAG}$
      $\;\longrightarrow\;$\\[0.4em]
      $\text{Sistema inteligente}$
      $\;\longrightarrow\;$
      $\textbf{Aplicación empresarial}$
    \end{block}
    \column{0.4\textwidth}
    \begin{exampleblock}{Lo que construiremos hoy}
      \begin{itemize}
        \item \textbf{RAG}: LLM con conocimiento del dominio
        \item \textbf{Redis}: LLM con memoria de sesión
        \item \textbf{LangChain}: LLM orquestado y trazable
        \item \textbf{JSON}: salida consumible por APIs
        \item \textbf{Pydantic}: salida validada y tipada
      \end{itemize}
    \end{exampleblock}
    \begin{block}{}
      \scriptsize
      \textbf{Reemplazabilidad}: una cadena LCEL bien diseñada
      permite cambiar \texttt{llama3.2} por \texttt{gpt-4o} sin
      modificar el resto del sistema.
    \end{block}
  \end{columns}
\end{frame}

% ═══════════════════════════════════════════════════════════════════
% DIAPOSITIVA 5: PROBLEMA DE NEGOCIO
% ═══════════════════════════════════════════════════════════════════
\begin{frame}{Caso de Estudio: Sistema de Minutas Alimentarias}
  \begin{columns}
    \column{0.55\textwidth}
    \begin{block}{Contexto del problema}
      \begin{itemize}
        \item Jardines infantiles JUNJI/VTF atienden niños de 0--3 años
        \item Cada niño posee un perfil de \textbf{alergias e intolerancias} documentado
        \item Se requieren minutas semanales por niño, por tiempo de comida
        \item Una nutricionista dedica $\approx$8 horas/semana a este proceso manual
        \item Errores en alergias tienen \textbf{consecuencias médicas graves}
      \end{itemize}
    \end{block}
    \vspace{0.3em}
    \begin{alertblock}{Restricciones del dominio}
      \begin{itemize}
        \item Normativa JUNJI circular N$^{\circ}$73 (2022): porciones, grupos alimentarios
        \item Tablas MINSAL de composición de alimentos
        \item Restricción calórica: 200--350 kcal por tiempo de comida
        \item Ingredientes disponibles en Chile con nombre local
      \end{itemize}
    \end{alertblock}
    \column{0.45\textwidth}
    \begin{exampleblock}{Salida esperada del sistema}
      \scriptsize
      El sistema debe retornar JSON estructurado:
      \vspace{0.4em}
      \begin{itemize}
        \item \texttt{nino\_id}: identificador del niño
        \item \texttt{semana}: semana ISO del año
        \item \texttt{dias[]}: array de 5 días hábiles
        \item \texttt{tiempos[]}: desayuno, almuerzo, once
        \item \texttt{ingredientes[]}: nombre + gramos
        \item \texttt{alergenos\_excluidos[]}: lista verificada
        \item \texttt{kcal}: valor calórico calculado
        \item \texttt{verificado\_alergenos}: booleano
      \end{itemize}
    \end{exampleblock}
    \vspace{0.3em}
    \begin{block}{Beneficio medible}
      \scriptsize
      De 8 horas/semana $\longrightarrow$ 15 minutos de revisión.
      El sistema genera; la nutricionista valida.
    \end{block}
  \end{columns}
\end{frame}

% ═══════════════════════════════════════════════════════════════════
% DIAPOSITIVA 6: ARQUITECTURA TRADICIONAL
% ═══════════════════════════════════════════════════════════════════
\begin{frame}{Arquitectura Tradicional: El Enfoque Incorrecto}

\begin{center}
\begin{tikzpicture}[
  box/.style={rectangle, rounded corners=4pt, minimum width=2.0cm,
              minimum height=0.75cm, text centered, font=\small,
              draw=black!50, fill=white, drop shadow},
  boxred/.style={rectangle, rounded corners=4pt, minimum width=2.0cm,
              minimum height=0.75cm, text centered, font=\small,
              draw=red!60, fill=red!8, drop shadow},
  arr/.style={-Stealth, thick, draw=black!55},
  arrd/.style={-Stealth, thick, draw=red!55, dashed},
]
  \node[box] (user) at (0,0)   {Usuario};
  \node[box] (app)  at (3,0)   {Aplicación\\Web};
  \node[boxred] (llm)  at (6,0)   {LLM\\(Ollama)};
  \node[boxred] (resp) at (9,0)   {Texto\\libre};

  \draw[arr] (user) -- (app);
  \draw[arr] (app) -- node[above, font=\scriptsize]{prompt sin contexto} (llm);
  \draw[arr] (llm)  -- (resp);

  % Problemas debajo
  \node[boxred, text width=2.2cm, align=center] (p1) at (2.5,-2.0)
        {\scriptsize\textbf{Sin memoria}\\\scriptsize sesión stateless};
  \node[boxred, text width=2.2cm, align=center] (p2) at (5.5,-2.0)
        {\scriptsize\textbf{Sin contexto}\\\scriptsize ignora normativa};
  \node[boxred, text width=2.2cm, align=center] (p3) at (8.5,-2.0)
        {\scriptsize\textbf{Sin estructura}\\\scriptsize texto no parseable};

  \draw[arrd] (app.south)  -- (p1.north);
  \draw[arrd] (llm.south)  -- (p2.north);
  \draw[arrd] (resp.south) -- (p3.north);
\end{tikzpicture}
\end{center}

\vspace{0.3em}
\begin{alertblock}{Por qué falla este enfoque}
  El LLM \textbf{no recuerda} consultas anteriores, \textbf{no conoce} la normativa JUNJI ni las tablas MINSAL,
  y sus respuestas en texto libre son \textbf{imposibles de integrar} automáticamente con un ERP o app.
\end{alertblock}
\end{frame}

% ═══════════════════════════════════════════════════════════════════
% DIAPOSITIVA 7: ANÁLISIS DEL FALLO
% ═══════════════════════════════════════════════════════════════════
\begin{frame}{Análisis Detallado del Fallo Arquitectural}
  \begin{columns}
    \column{0.33\textwidth}
    \begin{alertblock}{\centering Problema 1: Memoria}
      \begin{itemize}
        \item Cada consulta es un contexto \textbf{nuevo e independiente}
        \item El nutricionista debe re-ingresar alergias en cada sesión
        \item Sin historial, el LLM no aprende el perfil del niño
        \item \textbf{Escala}: 50 niños $\times$ 3 tiempos $\times$ 5 días = inmanejable
      \end{itemize}
    \end{alertblock}
    \column{0.33\textwidth}
    \begin{alertblock}{\centering Problema 2: Contexto}
      \begin{itemize}
        \item El LLM \textbf{no fue entrenado} con normativa JUNJI 2022
        \item Ignora tablas de composición MINSAL
        \item Las porciones son genéricas, no para 12--36 meses
        \item Puede \textbf{alucinar} ingredientes inexistentes en Chile
      \end{itemize}
    \end{alertblock}
    \column{0.33\textwidth}
    \begin{alertblock}{\centering Problema 3: Integración}
      \begin{itemize}
        \item Salida en lenguaje natural, no estructurada
        \item Imposible conectar con calendario, ERP o sistema de cocina
        \item Formato \textbf{inconsistente} entre ejecuciones
        \item Cualquier cambio en el prompt rompe el parseo
      \end{itemize}
    \end{alertblock}
  \end{columns}
  \vspace{0.5em}
  \begin{exampleblock}{Solución: tres capas complementarias}
    \centering
    \textbf{Memoria} (Redis) \quad + \quad \textbf{Contexto documental} (RAG) \quad + \quad \textbf{Salida estructurada} (JSON + Pydantic)
  \end{exampleblock}
\end{frame}

% ═══════════════════════════════════════════════════════════════════
% DIAPOSITIVA 8: ARQUITECTURA MODERNA
% ═══════════════════════════════════════════════════════════════════
\begin{frame}{Arquitectura Moderna Basada en IA}

\begin{center}
\begin{tikzpicture}[
  box/.style={rectangle, rounded corners=4pt, minimum width=1.9cm,
              minimum height=0.70cm, text centered, font=\scriptsize,
              draw=blue!55, fill=blue!8, drop shadow},
  boxg/.style={rectangle, rounded corners=4pt, minimum width=1.9cm,
              minimum height=0.70cm, text centered, font=\scriptsize,
              draw=green!60!black, fill=green!8, drop shadow},
  boxo/.style={rectangle, rounded corners=4pt, minimum width=1.9cm,
              minimum height=0.70cm, text centered, font=\scriptsize,
              draw=orange!70, fill=orange!10, drop shadow},
  boxp/.style={rectangle, rounded corners=4pt, minimum width=1.9cm,
              minimum height=0.70cm, text centered, font=\scriptsize,
              draw=purple!65, fill=purple!8, drop shadow},
  arr/.style={-Stealth, thick, draw=gray!65},
  arrd/.style={-Stealth, thick, draw=purple!55, dashed},
  lc/.style={rectangle, rounded corners=6pt, draw=blue!45, fill=blue!4, dashed,
             minimum width=7.5cm, minimum height=1.5cm},
]
  % LangChain container
  \node[lc] at (5.8, 0) {};
  \node[above, font=\scriptsize\bfseries\color{blue!65}] at (5.8, 0.75) {LangChain LCEL Chain};

  % Nodos dentro de la cadena
  \node[box] (prompt)  at (3.2, 0.0) {Prompt\\Template};
  \node[box] (ollama)  at (5.8, 0.0) {Ollama\\LLaMA 3.2};
  \node[boxg] (parser) at (8.4, 0.0) {JSON\\Parser};

  % Nodos externos
  \node[boxo] (user)   at (0.5, 0.0) {FastAPI\\(cliente)};
  \node[boxp] (rag)    at (3.2, -1.8) {RAG\\(Chroma)};
  \node[boxp] (redis)  at (5.8, -1.8) {Redis\\(historial)};
  \node[boxg] (output) at (11.2, 0.0) {JSON\\Minuta};

  % Flechas principales
  \draw[arr] (user)   -- (prompt);
  \draw[arr] (prompt) -- (ollama);
  \draw[arr] (ollama) -- (parser);
  \draw[arr, very thick, draw=green!60!black] (parser) -- (output);

  % Dependencias externas
  \draw[arrd] (rag)   -- node[right, font=\tiny]{contexto normativo} (prompt);
  \draw[arrd] (redis) -- node[right, font=\tiny]{sesión previa}      (ollama);
\end{tikzpicture}
\end{center}

\vspace{0.2em}
\begin{exampleblock}{Principio fundamental de diseño}
  El LLM es \textbf{uno de los componentes}, no el sistema completo.
  LangChain orquesta, RAG provee conocimiento, Redis provee memoria y
  el parser garantiza la estructura de salida.
\end{exampleblock}
\end{frame}

% ═══════════════════════════════════════════════════════════════════
% DIAPOSITIVA 9: RAG FUNDAMENTOS
% ═══════════════════════════════════════════════════════════════════
\begin{frame}{RAG: Retrieval-Augmented Generation}
  \begin{columns}
    \column{0.55\textwidth}
    \begin{block}{Definición formal}
      RAG es un patrón arquitectural que \textbf{aumenta} dinámicamente
      el contexto del prompt con fragmentos recuperados de una base de
      conocimiento externa, sin modificar los pesos del modelo.
    \end{block}
    \vspace{0.4em}
    \begin{block}{Aplicado a nuestro caso}
      \begin{itemize}
        \item \textbf{Corpus}: normativa JUNJI, tablas MINSAL, guías nutricionales pediátricas
        \item \textbf{Query}: ``Niño 2 años, alérgico a lactosa, almuerzo lunes''
        \item \textbf{Recuperación}: Top-4 fragmentos más similares semánticamente
        \item \textbf{Augmentation}: fragmentos inyectados en el prompt como contexto
      \end{itemize}
    \end{block}
    \column{0.45\textwidth}
    \begin{exampleblock}{RAG vs. Fine-tuning}
      \begin{tabular}{lll}
        \toprule
        \textbf{Criterio} & \textbf{RAG} & \textbf{Fine-tuning} \\
        \midrule
        Actualizable & Sí & No \\
        Trazable & Sí & No \\
        GPU entrenamiento & No & Sí \\
        Costo & Bajo & Alto \\
        Auditable & Sí & Difícil \\
        \bottomrule
      \end{tabular}
    \end{exampleblock}
    \vspace{0.4em}
    \begin{alertblock}{Limitación clave}
      RAG \textbf{no mejora el razonamiento} del LLM; amplía su
      base de conocimiento accesible en tiempo de inferencia.
    \end{alertblock}
  \end{columns}
\end{frame}

% ═══════════════════════════════════════════════════════════════════
% DIAPOSITIVA 10: RAG PIPELINE TIKZ
% ═══════════════════════════════════════════════════════════════════
\begin{frame}{RAG: Pipeline de Indexación y Recuperación}

\begin{center}
\begin{tikzpicture}[
  doc/.style={rectangle, rounded corners=2pt, minimum width=1.6cm,
              minimum height=0.58cm, text centered, font=\scriptsize,
              draw=orange!65, fill=orange!12},
  proc/.style={rectangle, rounded corners=2pt, minimum width=1.7cm,
               minimum height=0.58cm, text centered, font=\scriptsize,
               draw=blue!55, fill=blue!9},
  db/.style={cylinder, shape border rotate=90, aspect=0.35,
             minimum width=1.5cm, minimum height=0.72cm,
             text centered, font=\scriptsize,
             draw=purple!60, fill=purple!10},
  arr/.style={-Stealth, thick, draw=gray!65},
  arrd/.style={-Stealth, dashed, draw=purple!55},
  lbl/.style={font=\scriptsize\bfseries, text=gray!60},
]
  % Fase offline
  \node[lbl] at (0.8, 2.0) {OFFLINE: Indexación};
  \node[doc]  (pdf)    at (0.0, 1.3) {PDFs\\Normativa};
  \node[proc] (chunk)  at (2.2, 1.3) {Chunking\\512 tokens};
  \node[proc] (embed)  at (4.5, 1.3) {Encoder\\nomic-embed};
  \node[db]   (chroma) at (6.8, 1.3) {Chroma\\DB};

  \draw[arr] (pdf)   -- (chunk);
  \draw[arr] (chunk) -- (embed);
  \draw[arr] (embed) -- node[above, font=\tiny]{vectores float32} (chroma);

  % Fase online
  \node[lbl] at (0.8, 0.5) {ONLINE: Recuperación e inferencia};
  \node[doc]  (query)  at (0.0, -0.2) {Consulta\\usuario};
  \node[proc] (qemb)   at (2.2, -0.2) {Encoder\\(mismo modelo)};
  \node[proc] (sim)    at (4.5, -0.2) {Similitud\\coseno (MMR)};
  \node[proc] (topk)   at (6.8, -0.2) {Top-4\\fragmentos};
  \node[proc] (augp)   at (9.1, -0.2) {Prompt\\aumentado};
  \node[doc]  (llmout) at (11.4,-0.2) {LLM $\rightarrow$\\JSON};

  \draw[arr] (query) -- (qemb);
  \draw[arr] (qemb)  -- (sim);
  \draw[arr] (sim)   -- (topk);
  \draw[arr] (topk)  -- (augp);
  \draw[arr, very thick, draw=blue!65] (augp) -- (llmout);

  % Chroma conecta a similitud
  \draw[arrd] (chroma.south) -- node[right, font=\tiny]{búsqueda ANN} (sim.north);
\end{tikzpicture}
\end{center}

\vspace{0.2em}
\begin{block}{Invariante del embedding}
  El \textbf{mismo modelo de embedding} debe usarse en la fase de indexación
  y en la fase de recuperación. Para Ollama: \texttt{nomic-embed-text}
  (768 dims) o \texttt{mxbai-embed-large} (1024 dims).
\end{block}
\end{frame}

% ═══════════════════════════════════════════════════════════════════
% DIAPOSITIVA 11: REDIS CONCEPTO
% ═══════════════════════════════════════════════════════════════════
\begin{frame}{Redis: Memoria Contextual de Sesión}
  \begin{columns}
    \column{0.52\textwidth}
    \begin{block}{Por qué Redis para memoria de LLM}
      \begin{itemize}
        \item \textbf{In-memory}: latencia $<$ 1 ms (vs. PostgreSQL $\sim$5--20 ms)
        \item \textbf{TTL nativo}: sesiones expiran sin cron jobs
        \item \textbf{Listas}: historial de mensajes como estructura natural
        \item \textbf{Integración}: LangChain provee \texttt{RedisChatMessageHistory} nativo
        \item \textbf{Cluster}: escala horizontalmente sin cambios de código
      \end{itemize}
    \end{block}
    \vspace{0.3em}
    \begin{block}{Esquema de claves}
      \begin{itemize}
        \item Clave: \texttt{jardin:\{id\}:nino:\{id\}:sesion:\{fecha\}}
        \item Tipo: \texttt{LIST} de mensajes JSON serializados
        \item TTL: 86400 s (24 h) configurable
        \item Window: últimos N mensajes (presupuesto de tokens)
      \end{itemize}
    \end{block}
    \column{0.48\textwidth}
    \begin{exampleblock}{Tipos de memoria en LangChain}
      \begin{tabular}{ll}
        \toprule
        \textbf{Clase} & \textbf{Estrategia} \\
        \midrule
        \texttt{ConversationBuffer} & Historia completa \\
        \texttt{ConversationWindow} & Últimas N \\
        \texttt{ConversationSummary} & Resumen comprimido \\
        \texttt{ConversationEntity} & Entidades nombradas \\
        \texttt{VectorStore} & Semántica (RAG) \\
        \bottomrule
      \end{tabular}
    \end{exampleblock}
    \vspace{0.4em}
    \begin{alertblock}{Trade-off crítico}
      Más historial $=$ mejor contexto $=$ \textbf{más tokens} en el prompt
      $=$ mayor latencia y costo de inferencia.
      Usar ventana deslizante para controlar el tamaño.
    \end{alertblock}
  \end{columns}
\end{frame}

% ═══════════════════════════════════════════════════════════════════
% DIAPOSITIVA 12: REDIS CÓDIGO
% ═══════════════════════════════════════════════════════════════════
\begin{frame}[fragile]{Redis: Implementación con LangChain}
  \begin{columns}
    \column{0.56\textwidth}
\begin{lstlisting}[style=pythonstyle, title=Memoria con Redis y LCEL]
from langchain_community.chat_message_histories \
    import RedisChatMessageHistory
from langchain_core.runnables.history \
    import RunnableWithMessageHistory

REDIS_URL = "redis://localhost:6379"

def get_session_history(session_id: str):
    """Factory: un historial por sesion."""
    return RedisChatMessageHistory(
        session_id=session_id,
        url=REDIS_URL,
        ttl=86400,       # 24 horas en segundos
        key_prefix="minuta:",
    )

# Envolver la cadena base con memoria
chain_with_history = RunnableWithMessageHistory(
    base_chain,               # <- cadena LCEL
    get_session_history,      # <- factory Redis
    input_messages_key="input",
    history_messages_key="history",
)

# Primera consulta: informa alergias
r1 = chain_with_history.invoke(
    {"input": "Maria, 2 anios, alergia lactosa"},
    config={"configurable":
            {"session_id": "jard01:maria"}},
)
# Segunda consulta: Redis provee el contexto
r2 = chain_with_history.invoke(
    {"input": "Genera almuerzo del lunes"},
    config={"configurable":
            {"session_id": "jard01:maria"}},
)
\end{lstlisting}
    \column{0.44\textwidth}
    \begin{block}{Estructura en Redis}
      \scriptsize
      \textbf{Clave}: \texttt{minuta:jard01:maria}\\[0.3em]
      \textbf{Tipo}: \texttt{LIST} de strings JSON\\[0.4em]
      Cada elemento:
      \begin{lstlisting}[style=jsonstyle]
{
  "type": "human",
  "content": "Maria, 2 anios...",
  "additional_kwargs": {},
  "response_metadata": {}
}
      \end{lstlisting}
    \end{block}
    \vspace{0.3em}
    \begin{exampleblock}{Beneficio directo}
      En \texttt{r2}, el LLM ya sabe que María es
      alérgica a la lactosa \textbf{sin que se lo
      repita el usuario}. Redis lo inyecta
      automáticamente en el historial del prompt.
    \end{exampleblock}
  \end{columns}
\end{frame}

% ═══════════════════════════════════════════════════════════════════
% DIAPOSITIVA 13: LANGCHAIN CONCEPTO
% ═══════════════════════════════════════════════════════════════════
\begin{frame}{LangChain: Orquestación de Componentes LLM}
  \begin{columns}
    \column{0.55\textwidth}
    \begin{block}{LangChain Expression Language (LCEL)}
      LCEL define cadenas como \textbf{composición funcional} de
      componentes tipo \texttt{Runnable} mediante el operador \texttt{|}:
      \vspace{0.3em}
      \[
        \text{chain} = \underbrace{\text{prompt}}_{\text{Runnable}} \;|\;
                       \underbrace{\text{llm}}_{\text{Runnable}} \;|\;
                       \underbrace{\text{parser}}_{\text{Runnable}}
      \]
      \vspace{0.1em}
      Interfaz uniforme: \texttt{invoke()}, \texttt{stream()}, \texttt{batch()},
      \texttt{astream()}.
    \end{block}
    \vspace{0.3em}
    \begin{block}{Componentes clave para el caso de estudio}
      \begin{itemize}
        \item \texttt{ChatPromptTemplate}: plantilla con variables dinámicas
        \item \texttt{ChatOllama}: wrapper para Ollama local
        \item \texttt{JsonOutputParser}: garantiza dict Python en la salida
        \item \texttt{RunnablePassthrough.assign}: enriquece el diccionario de entrada
        \item \texttt{RunnableWithMessageHistory}: wrapper de memoria
      \end{itemize}
    \end{block}
    \column{0.45\textwidth}
    \begin{exampleblock}{Por qué no llamar Ollama directamente}
      \begin{itemize}
        \item \textbf{Retry}: reintentos automáticos ante fallos
        \item \textbf{Streaming}: tokens en tiempo real
        \item \textbf{Tracing}: LangSmith integrado (observabilidad)
        \item \textbf{Composición}: componentes intercambiables
        \item \textbf{Portabilidad}: cambiar LLM sin reescribir la lógica de negocio
        \item \textbf{Paralelismo}: \texttt{RunnableParallel}
      \end{itemize}
    \end{exampleblock}
    \vspace{0.3em}
    \begin{alertblock}{Versión usada en esta clase}
      \texttt{langchain $\geq$ 0.3.x} con LCEL nativo.
      La API antigua (\texttt{LLMChain}) está \textbf{deprecada}.
      No usar ejemplos de tutoriales anteriores a 2024.
    \end{alertblock}
  \end{columns}
\end{frame}

% ═══════════════════════════════════════════════════════════════════
% DIAPOSITIVA 14: LANGCHAIN CÓDIGO
% ═══════════════════════════════════════════════════════════════════
\begin{frame}[fragile]{LangChain: Construcción de la Cadena LCEL}
  \begin{columns}
    \column{0.56\textwidth}
\begin{lstlisting}[style=pythonstyle, title=chain\_builder.py -- Cadena LCEL completa]
from langchain_ollama import ChatOllama
from langchain_core.prompts import ChatPromptTemplate
from langchain_core.output_parsers \
    import JsonOutputParser
from langchain_core.runnables \
    import RunnablePassthrough
from operator import itemgetter

# 1. LLM local con salida JSON nativa
llm = ChatOllama(
    model="llama3.2",
    temperature=0.1,
    format="json",       # grammar-guided sampling
)

# 2. Plantilla: system + historial + human
prompt = ChatPromptTemplate.from_messages([
    ("system", "{system_prompt}"),
    ("placeholder", "{history}"),
    ("human",
     "Contexto normativo:\n{context}\n\n"
     "Solicitud: {input}"),
])

# 3. Cadena RAG + Prompt + LLM + Parser
base_chain = (
    RunnablePassthrough.assign(
        context=(
            itemgetter("input") | retriever
        )
    )
    | prompt
    | llm
    | JsonOutputParser()
)
\end{lstlisting}
    \column{0.44\textwidth}
    \begin{block}{Flujo de datos en la cadena}
      \scriptsize
      \begin{enumerate}
        \item \texttt{input}: consulta del nutricionista
        \item \texttt{assign(context=...)}: RAG recupera documentos
        \item \texttt{prompt}: ensambla system + history + context + human
        \item \texttt{llm}: Ollama genera tokens en formato JSON
        \item \texttt{JsonOutputParser}: convierte string $\rightarrow$ \texttt{dict}
      \end{enumerate}
    \end{block}
    \vspace{0.3em}
    \begin{exampleblock}{\texttt{format="json"} en Ollama}
      Instruye a \texttt{llama.cpp} (motor de inferencia)
      a usar \textbf{grammar-guided sampling}: solo puede
      generar tokens que mantengan el JSON sintácticamente
      válido. Garantía a nivel de motor, no de contenido.
    \end{exampleblock}
  \end{columns}
\end{frame}

% ═══════════════════════════════════════════════════════════════════
% DIAPOSITIVA 15: PRE-PROMPT ENGINEERING
% ═══════════════════════════════════════════════════════════════════
\begin{frame}[fragile]{Pre-Prompt Engineering: El Rol del Sistema}
  \begin{columns}
    \column{0.56\textwidth}
\begin{lstlisting}[style=pythonstyle, title=system\_prompt.py]
SYSTEM_PROMPT = """
Eres un sistema experto en nutricion infantil 
certificado segun normativa JUNJI/INTEGRA Chile.

REGLAS ABSOLUTAS (no negociables):
1. Responde UNICAMENTE en JSON valido.
2. NUNCA incluyas alimentos del listado de 
   alergenos del nino.
3. Las porciones son para ninos de 12-36 meses.
4. Cada tiempo de comida: 200-350 kcal.
5. Usa ingredientes disponibles en Chile 
   con nombre local (ej: "porotos", no "frijoles").
6. El campo "verificado_alergenos" debe ser true
   SOLO si revisaste explicitamente cada ingrediente.

ESQUEMA OBLIGATORIO:
{schema_json}

ANTE INCERTIDUMBRE:
Responde: {"error": "descripcion", 
           "alternativas": ["opcion1", "opcion2"]}

NUNCA inventes normativa ni valores nutricionales.
Cita el documento fuente si el contexto lo provee.
"""
\end{lstlisting}
    \column{0.44\textwidth}
    \begin{block}{Anatomía del system prompt}
      \begin{enumerate}
        \item \textbf{Rol}: identidad y certificación
        \item \textbf{Reglas absolutas}: restricciones inviolables numeradas
        \item \textbf{Dominio}: parámetros cuantitativos precisos
        \item \textbf{Esquema}: estructura JSON con variables
        \item \textbf{Fallback}: comportamiento explícito ante errores
        \item \textbf{Prohibiciones}: previene alucinaciones
      \end{enumerate}
    \end{block}
    \vspace{0.3em}
    \begin{alertblock}{Anti-patrón frecuente}
      \textit{``Responde con información nutricional para niños''} es
      una sugerencia, no una restricción. El LLM necesita
      \textbf{reglas explícitas y cuantificadas}, no intenciones vagas.
    \end{alertblock}
    \begin{exampleblock}{Técnica: few-shot en el prompt}
      Incluir 1--2 ejemplos JSON correctos aumenta
      la adherencia al esquema en $\approx$40\%.
    \end{exampleblock}
  \end{columns}
\end{frame}

% ═══════════════════════════════════════════════════════════════════
% DIAPOSITIVA 16: JSON ESTRUCTURADO
% ═══════════════════════════════════════════════════════════════════
\begin{frame}[fragile]{Respuestas Estructuradas: Diseño del Esquema JSON}
  \begin{columns}
    \column{0.52\textwidth}
\begin{lstlisting}[style=jsonstyle, title=Ejemplo de respuesta del sistema]
{
  "minuta": {
    "nino_id": "maria_jard01",
    "semana": "2026-W25",
    "alergenos_excluidos": ["lactosa"],
    "dias": [
      {
        "dia": "lunes",
        "tiempos": [
          {
            "nombre": "almuerzo",
            "plato": "Arroz con pollo al vapor",
            "ingredientes": [
              {"nombre": "arroz", "gramos": 60},
              {"nombre": "pechuga pollo", "gramos": 80},
              {"nombre": "zanahoria", "gramos": 30},
              {"nombre": "aceite oliva", "gramos": 5}
            ],
            "kcal": 292,
            "proteinas_g": 24.5,
            "carbohidratos_g": 38.1,
            "verificado_alergenos": true,
            "fuente_normativa": "JUNJI Circ. 73 Tab. 3"
          }
        ]
      }
    ]
  }
}
\end{lstlisting}
    \column{0.48\textwidth}
    \begin{block}{Validación con Pydantic v2}
\begin{lstlisting}[style=pythonstyle]
from pydantic import BaseModel, Field
from typing import List

class Ingrediente(BaseModel):
    nombre: str
    gramos: float = Field(gt=0, lt=500)

class TiempoComida(BaseModel):
    nombre: str
    plato: str
    ingredientes: List[Ingrediente]
    kcal: float = Field(gt=0, lt=600)
    verificado_alergenos: bool

class Minuta(BaseModel):
    nino_id: str
    semana: str
    alergenos_excluidos: List[str]
    dias: List[dict]
\end{lstlisting}
    \end{block}
    \vspace{0.2em}
    \begin{exampleblock}{Integración downstream}
      Este JSON es directamente consumible por:
      app móvil de padres, ERP de compras, sistema
      de cocina (tickets), calendario institucional.
    \end{exampleblock}
  \end{columns}
\end{frame}

% ═══════════════════════════════════════════════════════════════════
% DIAPOSITIVA 17: CONFIGURACIÓN BASE
% ═══════════════════════════════════════════════════════════════════
\begin{frame}[fragile]{Implementación: Configuración del Entorno}
  \begin{columns}
    \column{0.56\textwidth}
\begin{lstlisting}[style=pythonstyle, title=config.py -- Inicialización de componentes]
# requirements.txt
# langchain>=0.3.0          langchain-ollama>=0.2.0
# langchain-chroma>=0.1.0   langchain-community>=0.3.0
# redis>=5.0.0              pydantic>=2.0.0

from langchain_ollama import (
    ChatOllama, OllamaEmbeddings
)
from langchain_chroma import Chroma

LLM_MODEL   = "llama3.2"
EMBED_MODEL = "nomic-embed-text"
CHROMA_DIR  = "./chroma_db"
REDIS_URL   = "redis://localhost:6379"

# LLM: temperatura baja para respuestas consistentes
llm = ChatOllama(
    model=LLM_MODEL,
    temperature=0.1,
    format="json",
    num_ctx=4096,    # ventana de contexto
)

# Embeddings: mismo modelo en indexacion y consulta
embeddings = OllamaEmbeddings(model=EMBED_MODEL)

# Vector store persistente
vectorstore = Chroma(
    persist_directory=CHROMA_DIR,
    embedding_function=embeddings,
    collection_name="normativa_junji",
)
retriever = vectorstore.as_retriever(
    search_type="mmr",
    search_kwargs={"k": 4, "fetch_k": 20},
)
\end{lstlisting}
    \column{0.44\textwidth}
    \begin{block}{Modelos a descargar en Ollama}
\begin{lstlisting}[style=shellstyle]
# Terminal 1: servidor Ollama
$ ollama serve

# Terminal 2: descargar modelos
$ ollama pull llama3.2
$ ollama pull nomic-embed-text

# Verificar
$ ollama list
\end{lstlisting}
    \end{block}
    \vspace{0.3em}
    \begin{block}{Requisitos del sistema}
      \begin{itemize}
        \item \texttt{ollama} en \texttt{localhost:11434}
        \item \texttt{redis-server} en \texttt{localhost:6379}
        \item Python 3.11+, $\geq$ 8 GB RAM
        \item \texttt{llama3.2:3b-q4\_K\_M}: $\approx$ 2 GB VRAM
      \end{itemize}
    \end{block}
    \begin{exampleblock}{Para producción}
      Reemplazar \texttt{ChatOllama} por \texttt{ChatOpenAI}
      o \texttt{ChatAnthropic} sin tocar la cadena.
    \end{exampleblock}
  \end{columns}
\end{frame}

% ═══════════════════════════════════════════════════════════════════
% DIAPOSITIVA 18: RAG IMPLEMENTACIÓN
% ═══════════════════════════════════════════════════════════════════
\begin{frame}[fragile]{Implementación: Pipeline RAG con Chroma}
  \begin{columns}
    \column{0.56\textwidth}
\begin{lstlisting}[style=pythonstyle, title=rag\_indexer.py]
from langchain_community.document_loaders \
    import PyPDFLoader, DirectoryLoader
from langchain.text_splitter \
    import RecursiveCharacterTextSplitter

def build_vector_store(pdf_dir: str):
    """Indexa el corpus normativo completo."""
    loader = DirectoryLoader(
        pdf_dir, 
        glob="**/*.pdf",
        loader_cls=PyPDFLoader
    )
    docs = loader.load()
    print(f"Cargados {len(docs)} documentos")

    splitter = RecursiveCharacterTextSplitter(
        chunk_size=512,
        chunk_overlap=64,
        separators=["\n\n", "\n", ". ", " ", ""],
        length_function=len,
    )
    chunks = splitter.split_documents(docs)
    print(f"Generados {len(chunks)} fragmentos")

    vectorstore = Chroma.from_documents(
        documents=chunks,
        embedding=embeddings,
        persist_directory=CHROMA_DIR,
        collection_name="normativa_junji",
    )
    return vectorstore

# Ejecutar una vez al cambiar el corpus
if __name__ == "__main__":
    build_vector_store("./documentos")
\end{lstlisting}
    \column{0.44\textwidth}
    \begin{block}{Parámetros de chunking}
      \begin{tabular}{ll}
        \toprule
        \textbf{Parámetro} & \textbf{Valor} \\
        \midrule
        \texttt{chunk\_size} & 512 tokens \\
        \texttt{chunk\_overlap} & 64 tokens \\
        \texttt{retriever k} & 4 docs \\
        \texttt{fetch\_k} & 20 candidatos \\
        \texttt{search\_type} & MMR \\
        \bottomrule
      \end{tabular}
    \end{block}
    \vspace{0.3em}
    \begin{exampleblock}{MMR vs Similarity Search}
      \textbf{Similarity}: máxima relevancia individual.\\[0.2em]
      \textbf{MMR (Max Marginal Relevance)}: balancea relevancia
      y \textbf{diversidad} -- evita recuperar 4 fragmentos
      del mismo párrafo. Recomendado para corpus normativo.
    \end{exampleblock}
    \vspace{0.3em}
    \begin{alertblock}{Corpus mínimo recomendado}
      Circular JUNJI N$^{\circ}$73 (2022),
      Tablas MINSAL composición nutricional,
      Guías alimentarias MINSAL 0--3 años.
    \end{alertblock}
  \end{columns}
\end{frame}

% ═══════════════════════════════════════════════════════════════════
% DIAPOSITIVA 19: CADENA COMPLETA
% ═══════════════════════════════════════════════════════════════════
\begin{frame}[fragile]{Implementación: Cadena LangChain Completa}
  \begin{columns}
    \column{0.56\textwidth}
\begin{lstlisting}[style=pythonstyle, title=minuta\_chain.py -- Sistema completo]
from langchain_core.runnables \
    import RunnablePassthrough
from langchain_core.runnables.history \
    import RunnableWithMessageHistory
from operator import itemgetter

def build_minuta_chain(llm, retriever):
    prompt = ChatPromptTemplate.from_messages([
        ("system", SYSTEM_PROMPT),
        ("placeholder", "{history}"),
        ("human",
         "Contexto normativo:\n{context}\n\n"
         "Solicitud: {input}"),
    ])

    # Cadena RAG + LLM + Parser (sin memoria)
    base_chain = (
        RunnablePassthrough.assign(
            context=(
                itemgetter("input") | retriever
            )
        )
        | prompt
        | llm
        | JsonOutputParser()
    )

    # Envolver con memoria Redis
    return RunnableWithMessageHistory(
        base_chain,
        get_session_history,
        input_messages_key="input",
        history_messages_key="history",
    )

chain = build_minuta_chain(llm, retriever)

# Invocacion con sesion
result: dict = chain.invoke(
    {"input": "Almuerzo lunes para Maria, "
               "2 anios, alergia lactosa"},
    config={"configurable":
            {"session_id": "jard01:maria"}},
)
print(result["minuta"]["dias"][0])
\end{lstlisting}
    \column{0.44\textwidth}
    \begin{block}{Capas de la cadena final}
      \scriptsize
      \begin{enumerate}
        \item \textbf{RunnablePassthrough.assign}: enriquece el input
              con documentos RAG (Chroma)
        \item \textbf{Prompt}: ensambla system + historial Redis
              + contexto normativo + consulta
        \item \textbf{ChatOllama}: inferencia local llama3.2
              con grammar-guided JSON sampling
        \item \textbf{JsonOutputParser}: convierte string
              $\rightarrow$ \texttt{dict} Python tipado
        \item \textbf{RunnableWithMessageHistory}: persiste
              conversación en Redis y la inyecta en futuras invocaciones
      \end{enumerate}
    \end{block}
    \vspace{0.3em}
    \begin{exampleblock}{Resultado de \texttt{result}}
      \texttt{result} es un \texttt{dict} Python
      directamente serializable como \texttt{JSONResponse}
      en FastAPI o persistible en PostgreSQL.
    \end{exampleblock}
  \end{columns}
\end{frame}

% ═══════════════════════════════════════════════════════════════════
% DIAPOSITIVA 20: ACTIVIDAD 1
% ═══════════════════════════════════════════════════════════════════
\begin{frame}{Actividad 1: Primera Consulta al Sistema (30 min)}
  \begin{block}{Objetivo}
    Levantar el sistema mínimo viable e invocar la cadena para obtener
    la primera minuta estructurada en JSON.
  \end{block}
  \vspace{0.3em}
  \begin{columns}
    \column{0.5\textwidth}
    \begin{exampleblock}{Pasos a seguir}
      \begin{enumerate}
        \item Clonar el repositorio base de la clase
        \item Verificar Ollama activo: \texttt{ollama serve}
        \item Descargar modelos: \texttt{ollama pull llama3.2} y \texttt{nomic-embed-text}
        \item \texttt{pip install -r requirements.txt}
        \item Ejecutar \texttt{python rag\_indexer.py} con los PDFs de prueba
        \item Invocar la cadena y verificar el JSON de salida con \texttt{print(result)}
      \end{enumerate}
    \end{exampleblock}
    \column{0.5\textwidth}
    \begin{alertblock}{Criterios de aceptación}
      \begin{itemize}
        \item La cadena retorna un \texttt{dict} Python (no string)
        \item El JSON contiene \texttt{minuta.dias} con $\geq$ 1 día
        \item Todos los campos del esquema están presentes
        \item \texttt{verificado\_alergenos == true}
        \item \texttt{kcal} está en el rango 200--350
      \end{itemize}
    \end{alertblock}
    \vspace{0.4em}
    \begin{block}{Entregable}
      Captura del JSON generado para la consulta:\\
      \textit{``Niño 18 meses, sin alergias, almuerzo martes''}
    \end{block}
  \end{columns}
\end{frame}

% ═══════════════════════════════════════════════════════════════════
% DIAPOSITIVA 21: ACTIVIDAD 2
% ═══════════════════════════════════════════════════════════════════
\begin{frame}{Actividad 2: Memoria Contextual con Redis (25 min)}
  \begin{block}{Objetivo}
    Verificar empíricamente que el sistema mantiene contexto entre
    consultas mediante Redis sin que el usuario repita información.
  \end{block}
  \vspace{0.3em}
  \begin{columns}
    \column{0.5\textwidth}
    \begin{exampleblock}{Secuencia de consultas}
      \textbf{Consulta 1} (session\_id: \texttt{test:pedro}):
      \begin{quote}
        \scriptsize\itshape
        ``Pedro tiene 2 años y es alérgico a los mariscos
        y al huevo. Genera su desayuno del lunes.''
      \end{quote}
      \vspace{0.4em}
      \textbf{Consulta 2} (mismo session\_id, sin mencionar alergias):
      \begin{quote}
        \scriptsize\itshape
        ``Ahora genera su almuerzo del lunes.
        Respeta todas sus restricciones alimentarias.''
      \end{quote}
      \vspace{0.4em}
      \textbf{Consulta 3} (session\_id \textbf{diferente}: \texttt{test:pedro2}):
      \begin{quote}
        \scriptsize\itshape
        ``Genera el almuerzo del lunes.''
      \end{quote}
    \end{exampleblock}
    \column{0.5\textwidth}
    \begin{alertblock}{¿Qué verificar?}
      \begin{enumerate}
        \item Consulta 2: el almuerzo NO contiene mariscos ni huevo
        \item Verificar en Redis:\\
              \texttt{redis-cli LRANGE test:pedro 0 -1}
        \item Consulta 3 (nueva sesión): el LLM pide más contexto o genera minuta genérica
      \end{enumerate}
    \end{alertblock}
    \vspace{0.3em}
    \begin{block}{Pregunta de análisis}
      ¿Por qué la consulta 3 falla o produce resultados distintos?
      ¿Qué implica esto para el diseño del sistema en producción?
      Documenta la diferencia de comportamiento.
    \end{block}
  \end{columns}
\end{frame}

% ═══════════════════════════════════════════════════════════════════
% DIAPOSITIVA 22: ACTIVIDAD 3
% ═══════════════════════════════════════════════════════════════════
\begin{frame}[fragile]{Actividad 3: Minuta Semanal con Validación Pydantic (40 min)}
  \begin{block}{Objetivo}
    Extender la cadena para generar una minuta completa de 5 días
    con validación estricta mediante Pydantic v2.
  \end{block}
  \vspace{0.3em}
  \begin{columns}
    \column{0.5\textwidth}
    \begin{exampleblock}{Tareas de implementación}
      \begin{enumerate}
        \item Crear los modelos Pydantic completos: \texttt{Ingrediente}, \texttt{TiempoComida}, \texttt{DiaSemana}, \texttt{MinutaSemanal}
        \item Reemplazar \texttt{JsonOutputParser} por \texttt{PydanticOutputParser}
        \item Inyectar el formato Pydantic en el system prompt con \texttt{parser.get\_format\_instructions()}
        \item Implementar \textbf{retry automático}: si \texttt{ValidationError}, reenviar el error al LLM para autocorrección
        \item Persistir la minuta validada en Redis con TTL de 7 días
      \end{enumerate}
    \end{exampleblock}
    \column{0.5\textwidth}
    \begin{block}{Hint: OutputFixingParser}
\begin{lstlisting}[style=pythonstyle]
from langchain.output_parsers \
    import OutputFixingParser

fixing_parser = OutputFixingParser\
    .from_llm(
        parser=pydantic_parser,
        llm=llm,
        max_retries=3,
    )
\end{lstlisting}
    \end{block}
    \vspace{0.2em}
    \begin{alertblock}{Desafío avanzado (opcional)}
      Implementar \textbf{doble verificación}: después de obtener el JSON,
      usar un segundo LLM call para auditar que ningún ingrediente contenga
      el alérgeno declarado.\\[0.3em]
      Reflexiona: ¿es esto RAG, cadena o agente?
    \end{alertblock}
  \end{columns}
\end{frame}

% ═══════════════════════════════════════════════════════════════════
% DIAPOSITIVA 23: PREGUNTAS DE REFLEXIÓN
% ═══════════════════════════════════════════════════════════════════
\begin{frame}{Preguntas de Reflexión}
  \begin{columns}
    \column{0.5\textwidth}
    \begin{block}{Diseño arquitectural}
      \begin{enumerate}
        \item Si la normativa JUNJI se actualiza en marzo 2027,
              ¿qué componente del sistema debe modificarse y cómo
              se minimiza el impacto en los demás?
        \vspace{0.35em}
        \item ¿Es suficiente \texttt{format=``json''} en Ollama para
              garantizar JSON válido en producción? ¿Qué casos de borde
              pueden romper esto incluso con grammar-guided sampling?
        \vspace{0.35em}
        \item ¿Cómo diseñarías el sistema para 500 niños simultáneos?
              ¿Qué componentes escalan horizontalmente y cuáles son
              el cuello de botella?
      \end{enumerate}
    \end{block}
    \column{0.5\textwidth}
    \begin{exampleblock}{Evaluación crítica}
      \begin{enumerate}
        \setcounter{enumi}{3}
        \item El nutricionista reporta: \textit{``El sistema generó una
              receta que cumple el esquema JSON pero usa palmito, que
              no está en el proveedor del jardín''}. ¿Qué capa del
              sistema falló? ¿Cómo se previene arquitecturalmente?
        \vspace{0.35em}
        \item ¿Cuál es la diferencia fundamental entre esta cadena
              y un \textbf{agente LLM}? ¿Qué necesitaría agregarse
              para convertirla en un agente reactivo con herramientas?
        \vspace{0.35em}
        \item ¿Es ético usar IA para tomar decisiones sobre la
              alimentación de niños menores de 3 años sin supervisión
              humana en el loop de aprobación?
      \end{enumerate}
    \end{exampleblock}
  \end{columns}
\end{frame}

% ═══════════════════════════════════════════════════════════════════
% DIAPOSITIVA 24: ARQUITECTURA FINAL TikZ
% ═══════════════════════════════════════════════════════════════════
\begin{frame}{Arquitectura Final Integrada del Sistema}

\begin{center}
\begin{tikzpicture}[
  b/.style={rectangle, rounded corners=3pt, minimum width=1.65cm,
            minimum height=0.52cm, text centered,
            font=\fontsize{5.8}{7}\selectfont,
            draw=blue!55, fill=blue!8},
  bg/.style={rectangle, rounded corners=3pt, minimum width=1.65cm,
             minimum height=0.52cm, text centered,
             font=\fontsize{5.8}{7}\selectfont,
             draw=green!55!black, fill=green!8},
  bo/.style={rectangle, rounded corners=3pt, minimum width=1.65cm,
             minimum height=0.52cm, text centered,
             font=\fontsize{5.8}{7}\selectfont,
             draw=orange!65, fill=orange!10},
  bp/.style={rectangle, rounded corners=3pt, minimum width=1.65cm,
             minimum height=0.52cm, text centered,
             font=\fontsize{5.8}{7}\selectfont,
             draw=purple!60, fill=purple!9},
  br/.style={rectangle, rounded corners=3pt, minimum width=1.65cm,
             minimum height=0.52cm, text centered,
             font=\fontsize{5.8}{7}\selectfont,
             draw=red!55, fill=red!8},
  arr/.style={-Stealth, draw=gray!60, thick},
  arrd/.style={-Stealth, draw=purple!55, dashed, thick},
  grp/.style={rectangle, rounded corners=5pt,
              draw=gray!40, fill=gray!3, dashed},
  lbl/.style={font=\fontsize{5.5}{6.5}\selectfont\bfseries, text=gray!60},
]

  % ── Capa corpus ────────────────────────────────────────────────
  \node[lbl] at (0.6, 3.15) {CORPUS NORMATIVO};
  \node[bo] (pdf1)  at (-0.1, 2.55) {JUNJI\\Circ. 73};
  \node[bo] (pdf2)  at ( 1.7, 2.55) {MINSAL\\Tablas};
  \node[b]  (chunk) at ( 3.5, 2.55) {Chunking\\512 tok};
  \node[b]  (emb)   at ( 5.4, 2.55) {Embeddings\\nomic-embed};
  \node[bp] (chroma)at ( 7.3, 2.55) {Chroma DB\\vectores};

  \draw[arr] (pdf1)  -- (chunk);
  \draw[arr] (pdf2)  -- (chunk);
  \draw[arr] (chunk) -- (emb);
  \draw[arr] (emb)   -- (chroma);

  % ── Capa cliente ───────────────────────────────────────────────
  \node[lbl] at (0.3, 1.65) {CLIENTE};
  \node[bg] (api)    at (-0.1, 1.05) {FastAPI\\POST /minuta};
  \node[bg] (schema) at ( 1.7, 1.05) {Pydantic\\validación};

  \draw[arr] (api) -- (schema);

  % ── LangChain chain ────────────────────────────────────────────
  \begin{scope}[on background layer]
    \node[grp, fit={(3.0,0.62)(9.6,1.52)}] {};
  \end{scope}
  \node[lbl] at (4.3, 1.52) {LANGCHAIN LCEL CHAIN};

  \node[b]  (rpass)  at ( 3.5, 1.05) {Runnable\\Passthrough};
  \node[b]  (prompt) at ( 5.4, 1.05) {Prompt\\Template};
  \node[br] (ollama) at ( 7.3, 1.05) {Ollama\\LLaMA 3.2};
  \node[bg] (jparse) at ( 9.1, 1.05) {JSON\\Parser};

  \draw[arr] (schema) -- (rpass);
  \draw[arr] (rpass)  -- (prompt);
  \draw[arr] (prompt) -- (ollama);
  \draw[arr] (ollama) -- (jparse);

  % ── Memoria ────────────────────────────────────────────────────
  \node[lbl] at (0.3, 0.1) {MEMORIA};
  \node[bp] (redis) at (1.7, -0.42) {Redis\\SessionHistory};

  \draw[arrd] (redis.east) to[bend right=12]
              node[below, font=\fontsize{4.5}{5}\selectfont]{historial sesión}
              (prompt.south);
  \draw[arrd] (api.south)  --
              node[left, font=\fontsize{4.5}{5}\selectfont]{session\_id}
              (redis.north);

  % ── RAG retriever ──────────────────────────────────────────────
  \draw[arrd, purple!60] (chroma.south) to[bend left=8]
       node[right, font=\fontsize{4.5}{5}\selectfont]{top-4 docs}
       (rpass.north);

  % ── Output ─────────────────────────────────────────────────────
  \node[lbl] at (10.6, 1.65) {OUTPUT};
  \node[bg] (out)  at (10.9, 1.05) {JSON\\Minuta};
  \node[bg] (pyd2) at (10.9, 0.25) {Pydantic\\v2 check};
  \node[bo] (db)   at (10.9,-0.42) {PostgreSQL\\/ Redis TTL7d};

  \draw[arr, very thick, draw=green!55!black] (jparse) -- (out);
  \draw[arr] (out)  -- (pyd2);
  \draw[arr] (pyd2) -- (db);

\end{tikzpicture}
\end{center}

\end{frame}

% ═══════════════════════════════════════════════════════════════════
% DIAPOSITIVA 25: CONCLUSIONES
% ═══════════════════════════════════════════════════════════════════
\begin{frame}{Conclusiones}
  \begin{columns}
    \column{0.5\textwidth}
    \begin{block}{Lo que construimos hoy}
      \begin{itemize}
        \item Un sistema de software real donde el \textbf{LLM es un componente}, no la aplicación completa
        \item Pipeline RAG con corpus normativo actualizable sin reentrenar el modelo
        \item Memoria de sesión persistente y con expiración automática via Redis
        \item Salidas JSON validadas con Pydantic, \textbf{directamente consumibles} por APIs externas
        \item Una cadena LCEL reemplazable (cambiar LLM sin reescribir la lógica)
      \end{itemize}
    \end{block}
    \vspace{0.3em}
    \begin{exampleblock}{La narrativa completada}
      \centering
      \textbf{LLM solo}
      $\;\longrightarrow\;$
      \textbf{LLM + RAG}
      $\;\longrightarrow\;$
      \textbf{Sistema inteligente}
      $\;\longrightarrow\;$
      \textbf{Aplicación empresarial}
    \end{exampleblock}
    \column{0.5\textwidth}
    \begin{block}{Próximas clases}
      \begin{itemize}
        \item \textbf{Clase 10}: Agentes LLM -- de cadenas a agentes reactivos con LangGraph
        \item \textbf{Clase 11}: Evaluación -- LLM-as-Judge, RAGAS, benchmarks de RAG
        \item \textbf{Clase 12}: Observabilidad -- LangSmith, trazas en producción, costos
        \item \textbf{Clase 13}: Seguridad -- prompt injection, PII, guardrails y jailbreaking
      \end{itemize}
    \end{block}
    \vspace{0.3em}
    \begin{alertblock}{El principio central de esta clase}
      Un LLM sin contexto de dominio es un asistente genérico.
      Un LLM con RAG + memoria + restricciones estructurales
      es un \textbf{sistema inteligente listo para producción}.
    \end{alertblock}
  \end{columns}
\end{frame}

\end{document}
